\clearpage
\section*{2026-09-16: Document review supports Sullivan's full earlier parcel}
\textit{Agent-drafted research entry.}

The first documentary batch reviews the eleven post-policy parents with an
exact-99 constituent and unresolved land in the full-sample screen. Together
these parents contain 23 constituent filings, twenty exact-99 constituents,
and 2,207 units. Three GPT-5.6 Sol agents researched bounded groups, with
supervisory source inspection and review before accepting a finding.

Sullivan/Empire's four 99-unit filings are supported by the complete earlier
38,118-square-foot parcel. The supervisor inspected and saved the five-page
ACRIS zoning declaration, document 2026030900544002, dated March 2, 2026 and
recorded March 9. Page 3 explicitly partitions old Brooklyn block 1306 lot
28 into tentative lots 20, 25, 28, and 30 and names all four constituent
addresses. Page 5 gives the same external boundary. The diagram contains no
fifth remainder. Its approximate dimensional area is 38,119 square feet,
consistent with the recorded 38,118 after rounding; the administrative area
remains the adopted measure. The former lot was assembled in 2021, before
the 2023 reference map. Missing completed successor polygons therefore leave
a source-coverage gap that the recorded diagram resolves.

A small manual source table records this complete-parcel crosswalk. The audit
checks it against the saved reference parcel identifiers and recorded area,
while retaining independent filing and shared-land flags. This reduces the
unresolved audit count from 160 to 159 of 877 parents: 119 historical and
forty post-policy. Ten of the eleven exact-99 parents remain unresolved.
Sullivan's production land value already equals 38,118; this finding changes
its audit classification, not its production area or its units.

The review also identifies a possible missed economic-parent link. The
99-unit filing at 3008 Godwin Terrace and the 65-unit filing at 220 Kimberly
Place are separate parents in the current panel. Both were filed March 31,
2025 with the same owner and architect, and both point to old lot 78 in the
parcel history. The link rule does not accept their shared APPBBL history
because its recorded change is after filing; neither job supplies an explicit
cross-reference or common project code. A May 2025 recorded agreement's
inspected opening pages distinguish developer lot 78 from neighboring owner
lot 99 within a shared zoning lot. This supports further investigation but
does not yet establish the allocation among the two buildings, retained land,
and other uses. The proposed 164-unit parent is not adopted.

The 165th Street residential-unit contract names lots 65 and 85 and parts
of lots 30 and 89, without dimensions; the eight-parcel assemblage remains
unallocated. Beach 30th's filed City Planning plan describes a
63,346.60-square-foot zoning lot for two buildings that also includes parcels
with other live filings and omits the third parent building. These documents
clarify why the earlier economic-parent boundary remains unresolved. Later
polygon areas were retained as comparison evidence. Two initial agent proposals to resolve
land from those measurements were rejected. A Broadway zoning drawing was
available only through a search-index extract; its scan could not be retrieved
and its claimed contents were not treated as verified evidence. HDB Class A
unit counts retain priority over differently defined DOB totals. All parents
remain in the production panels, with membership, dates, units, land values,
and weights unchanged.

\textit{Reproduction.} Working changes based on \texttt{a665b53}; run
\texttt{make dof-site-review}. The existing audit reads
\path{parent_site_document_review.csv}. The source scan, manual inspection
notes, and frozen public API extracts are preserved under
\path{tasks/audits/download_parent_review_documents/code/} in the dated
\texttt{site\_research} folders. The supervisor's report,
\path{site_research_supervised.md}, separates the accepted finding from agent
proposals and unverified leads. A legal zoning lot does not by itself define
the economic parent or the wage-assessment unit.
